\section{语义规约层：失效模式分类（C-1）}\label{sec:failure}

\subsection{为什么需要分类}

若把 L2 的失败统称为“规约没做好”，改进方向只能是“换个更强的模型”，这既不可诊断也不可归因。本节把失败细化为四类。分类的判据是\emph{可操作性}：给出一条失败样例，应当能通过检查确定它属于哪一类，从而知道该修哪里。我们明确不主张该分类是完备的。

\begin{table}[t]
\centering
\small
\caption{语义规约层的四类失效模式。$D$ 列为“该失效是否会被纯答案级验证 V0 检出”。}
\label{tab:failures}
\begin{tabularx}{\textwidth}{@{}p{2.6cm} p{3.5cm} X c@{}}
\toprule
\textbf{失效模式} & \textbf{触发条件} & \textbf{表现} & \textbf{V0 可检} \\
\midrule
\textbf{F1 歧义未消解} & 自然语言表述多解 & 规约时静默选定一种解释，未把该选择记入假设集 $H$，后续验证与结论都建立在未声明的解释上 & 否 \\
\addlinespace
\textbf{F2 类型不完整} & 变量域/量词范围/定义域缺失 & $\sigma$ 中变量的类型或定义域 $D_v$ 丢失，使约束检查无法判定合法性，或把只在子域成立的结论外推到全域 & 否 \\
\addlinespace
\textbf{F3 不可逆} & $S$ 是有损映射 & 形式表示无法反译回可读的自然语言陈述，验证判定无法归因到原问题的哪一步 & 否 \\
\addlinespace
\textbf{F4 语义漂移} & 规约正确但后续改写失守 & 化简/代入/近似步骤脱离了 $H$ 中记录的假设，形式上仍通过检查，语义上已不是原问题 & 否 \\
\bottomrule
\end{tabularx}
\end{table}

\subsection{四类失效模式}

\paragraph{F1 歧义未消解（Ambiguity）。}自然语言数学问题的表述常允许多种合理读法，例如“$\triangle ABC$ 中的中线”既可指某条中线，也可指中线族；“随机选取”既可指均匀也可指其他分布。人在解题时会根据语境消解歧义并把它作为隐含前提。规约层若静默选择一种读法而不写入 $H$，则 $S$ 的输出 $\sigma$ 在形式上是自洽的，但它的\emph{语义锚点}缺失。危害在于：这类失败不会表现为任何检查报错，只会在与出题人意图不符时暴露，而那时错误已经过完整条链条。

\paragraph{F2 类型不完整（Type Incompleteness）。}$\sigma = \langle V,H,G_0,\tau\rangle$ 中的变量必须携带类型 $T_v$ 与定义域 $D_v$。规约时常犯的错误是只保留变量的“名字与数值”，丢掉其类型与定义域约束。典型后果是把仅在 $n \ge 3$ 成立的组合恒等式外推到 $n = 2$，或把实数域的运算规则用于整数环上的模运算。此类的可检出性依赖于检查器是否被给予类型信息——这正是 V2 级验证的动机。

\paragraph{F3 不可逆（Non-invertibility）。}若规约结果无法回译为自然语言陈述，则 L3 的判定只能针对 $\sigma$ 成立，无法传递回原问题（推论 R3）。形式上，这要求存在 $S^{-1}$ 使 $S^{-1}(S(r))$ 与 $r$ 语义一致。实践中不可逆的常见来源是：把多个自然语言步骤压成一条符号等式（丢失中间解释）、或把文字条件转换成无名约束（丢失可读的对应关系）。

\paragraph{F4 语义漂移（Semantic Drift）。}这一类与前三类不同：规约本身是正确的，问题出在规约之后的处理。一旦 $\sigma$ 离开 $H$ 中的假设（例如在后续化简中默认某量为正、忽略边界条件），后续验证仍会对\emph{已改变的问题}给出一致判定。语义漂移之所以危险，在于它同时满足“检查通过”与“结论错误”两个条件，且无法被 V0 或纯答案比对捕获。

\subsection{与既有分析的接口}

上述分类与文献中的若干已知现象可以对齐：F4 与“自我纠错把正确答案改错”的现象同源\citep{huang2023large}——纠错步骤脱离了原假设；F1/F2 解释了为何仅改变表面表述会显著降低准确率\citep{mirzadeh2024gsm}，因为模型重现的是训练分布中的表面模式而非结构；F3 正是自动形式化路线中“翻译正确性与覆盖性同时受限”在非形式化场景下的对应物\citep{szegedy2020promising,wu2022autoformalization}。分类的价值不在于解释一切，而在于把上述现象归入统一的可修复结构。
